\centerline{\bf \Huge Правильный перебор II}

\begin{flushright}
        \scriptsize{Здесь может быть Ваша цитата}
\end{flushright}

\problem{\textbf{0}} Начальник обратил внимание, что дата проведения третьего дня Летних математических сборов, записанная восемью цифрами (16.07.2025) обладает интересной особенностью: сумма первых четырёх цифр на 5 больше, суммы цифр номера года. А какие ещё даты в 2025 году имеют такое же свойство?

\begin{flushright} \textbf{2 балла}\\ 
\end{flushright}

\problem{\textbf{1}} Перечислите все четвёрки натуральных чисел, дающих в сумме
15.

\begin{flushright} \textbf{3 балла}\\ 
\end{flushright}


\problem{\textbf{2}} В январе некоторого года было 4 понедельника и 4 пятницы.
Каким днем недели могло быть 20-е число этого месяца? (2 попытки!)



\begin{flushright} \textbf{3 балла}\\ 
\end{flushright}

\problem{\textbf{3}} На гранях кубика расставлены числа от 1 до 6.
Кубик бросили два раза. В первый раз сумма чисел на четырёх боковых гранях оказалась
равна 12, во второй — 15. Какое число написано на грани, противоположной той, где написана
цифра 3? (2 попытки!)




\begin{flushright} \textbf{3 балла}\\ 
\end{flushright}


\problem{\textbf{4}}  Некоторое четырёхзначное число является точным
квадратом. Если убрать первую цифру слева, то оно станет точным кубом, а если убрать две
первые цифры, то оно станет четвёртой степенью целого числа. Найдите это число.



\begin{flushright} \textbf{3 балла}\\ 
\end{flushright}

\problem{\textbf{5}}  На конкурсе «А ну-ка, чудища!» стоят в ряд 15 драконов. У соседей число голов отличается на 1. Если у дракона больше голов, чем у обоих его
соседей, его считают хитрым, если меньше, чем у обоих соседей, — сильным, остальных (в том
числе стоящих с краю) считают обычными. В ряду есть ровно четыре хитрых дракона — с 4,
6, 7 и 7 головами и ровно три сильных — с 3, 3 и 6 головами. У первого и последнего драконов
голов поровну. Приведите хотя бы два примера того, как такое могло быть.


\begin{flushright} \textbf{3 балла}\\ 
\end{flushright}